% ===== Opening: title, hook, the book and the series, today's plan (~5 min) =====

\webinartitle{3}{Գլուխ 3. Վախ. վախենալու բան չկա}%
  {\chip{clover}{սնահավատություն} \chip{ufo}{դավադրություններ} \chip{spider}{ֆոբիաներ} \chip{pumpkin}{սարսափ ֆիլմեր} \chip{thumbsdown}{քննադատություն}}%
  {scared}{Վտանգ է հայտնաբերվել։ Ամենուր։}{{bat}{blackcat}{snake}}

% Hook: the chapter's own opening list of worries, as background processes
\begin{frame}{Ի՞նչն է քեզ հիմա անհանգստացնում}
  \begin{columns}[c]
    \begin{column}{0.64\textwidth}
      \begin{popup}{ուղեղ.exe}
        {\blackfont Ֆոնային անհանգստություններ}
        \begin{itemize}
          \item<2-> Կկարողանա՞մ գազի հաշիվը վճարել
          \item<3-> Մայրիկը վաղուց չի զանգել. լա՞վ է
          \item<4-> Միսը մի շաբաթ է՝ սառնարանում. փչացե՞լ է
          \item<5-> 9 տարեկանում տաբատս ընկավ. դեռ հիշու՞մ են
          \item<6-> Էս ի՞նչ ձայն էր. առնե՞տ է։ Ժանտախտով չլինի՞
        \end{itemize}
        \vspace{-1mm}
      \end{popup}
    \end{column}
    \begin{column}{0.3\textwidth}
      \centering
      \mascot[2.4cm]{panic}\par\medskip
      \onslide<7->{%
      \begin{chatbox}
        Գրիր մի բան, որի համար այսօր անհանգստացել ես։
      \end{chatbox}}
    \end{column}
  \end{columns}
\end{frame}

% The book and where this chapter sits in the series (for newcomers)
\begin{frame}{Գիրքը, հեղինակը և մեր շարքը}
  \begin{columns}[T]
    \begin{column}{0.5\textwidth}
      \begin{card}
        {\blackfont\Large The Idiot Brain} {\small (2016)}\par
        «Ապուշ ուղեղը»\par\medskip
        \begin{itemize}
          \item \textbf{Դին Բըրնեթ}՝ նեյրոգիտնական, Քարդիֆի համալսարան
          \item Guardian-ի «Brain Flapping» բլոգը, ստենդափ
          \item Գրված է կատակով, բայց գիտական աղբյուրներով
        \end{itemize}
        \vspace{-1mm}
      \end{card}
    \end{column}
    \begin{column}{0.44\textwidth}
      {\blackfont\small Գրքի 8 գլուխը}\par\smallskip
      {\color{plum!55!cream}\small \emoji{check}\ Ներածություն, 1, 2՝ վեբինար 1-2}\par\smallskip
      \chip[brainpink]{scream}{3. Վախ. այսօր}\par\smallskip
      {\color{plum!55!cream}\small
      4. Խելք \quad 5. Ընկալում\\
      6. Անհատականություն \quad 7. Ուրիշ մարդիկ\\
      8. Երբ ուղեղը «փչանում» է}\par\smallskip
      \begin{card}[colback=mint!35!cardwhite]
        Նախորդներին չե՞ս եղել։ Ոչինչ. ամեն մեկն ինքնուրույն է։
      \end{card}
    \end{column}
  \end{columns}
\end{frame}

% Plan for today: the four sections of chapter 3
\begin{frame}{Այսօր}
  \begin{columns}[T]
    \begin{column}{0.48\textwidth}
      \begin{card}[colback=lavender!40!cream]
        \emoji[1.3em]{clover}\ {\blackfont Սնահավատություն ու դավադրություններ}\par\smallskip
        Ինչու ենք կապեր տեսնում այնտեղ, որտեղ դրանք չկան
      \end{card}
      \begin{card}[colback=peach!45!cream]
        \emoji[1.3em]{pumpkin}\ {\blackfont Վախը՝ որպես հաճույք}\par\smallskip
        Ինչու ենք կամովին գնում սարսափ ֆիլմի ու ցատկում պարաշյուտով
      \end{card}
    \end{column}
    \begin{column}{0.48\textwidth}
      \begin{card}[colback=mint!40!cream]
        \emoji[1.3em]{microphone}\ {\blackfont Ֆոբիաներ ու կարաոկե}\par\smallskip
        Ինչու ենք սարսափում բաներից, որոնք գիտենք, որ անվտանգ են
      \end{card}
      \begin{card}[colback=brainpink!35!cream]
        \emoji[1.3em]{thumbsdown}\ {\blackfont Քննադատություն}\par\smallskip
        Ինչու է մեկ վատ խոսքը ավելի ուժեղ, քան տասը լավը
      \end{card}
    \end{column}
  \end{columns}
  \vspace{2mm}
  \centering
  {\small \emoji{coffee}\ Մեջտեղում՝ 5 րոպե ընդմիջում \qquad \emoji{speech}\ Հարցերը գրիր չաթում ցանկացած պահի}
\end{frame}
